\scanpage{158}
\mspage{16}
% archive sheet number: 158
\uline{NB} Si $X$ loc. fbt filtrante à g., \deleted{tout} $X'$ étale sur $X$ itou.
\uline{NB} $X$ loc. fbt filtrante à g $\Longleftrightarrow$ les compos. \unsure{ses} connexes \deleted{sont} fbt filtr. à g., ou encore si $X$ \unsure{est} 0-connexe, $X$ loc. fbt filtr. à g $\Longleftrightarrow X$ est \inserted{\unsure{pseudo-filtrante}} fbt. filtr. à g.
\margin{$\longleftarrow$ \uuline{Faux}}
On a en effet $\Longleftarrow$ \unsure{tautologique}*, et $\Longrightarrow$ \deleted{\unsure{pour un}} en montrant \deleted{sur la grande \unsure{cat} de \unsure{dire} puisque $a$ et $b$} \inserted{\unsure{à PS1 abs} (sauf ?) alors} que \deleted{\unsure{le}} $R(a,b)$ défini par $X_{/[a]}\cap X_{/[b]}\ne\varnothing$
% slanted
\margin{\uuline{faux}, ainsi
\diagram{\begin{tikzcd}[ampersand replacement=\&,column sep=small]
 0\arrow[r,bend left,"\alpha"]\arrow[r,bend right,"\beta"'] \&1
\end{tikzcd}}
est 0-conn. et fbt. filtr. à g., mais non pas loc. fbt. filtr. à g. ; car
\diagram{\begin{tikzcd}[ampersand replacement=\&,row sep=0pt,column sep=small]
 0\arrow[dr,"\alpha"] \&{}\\{} \&1\\0\arrow[ur,"\beta"'] \&{}
\end{tikzcd}}
ne se complète pas en carré \uline{commutatif}}
est une relation d'équivalence, impliquée par la relation $\deleted{\unsure{S}}\unsure{C|S}(a,b)$ : \{ $\mathrm{Hom}(a,b)\ne\varnothing$ \uline{ou} $\mathrm{Hom}(b,a)\ne\varnothing$\} donc \deleted{\unsure{est plus}} la relation d'équivalence (\unsure{visiblement plus|d'isomotopie}) engendrée par $C(a,b)$ est plus fine que $R(a,b)$.

c) On dit que $X$ est \uline{filtrante à g.}, \inserted{ou \uline{cofiltrante},} si elle est (\unsure{fbt} filtrante à g. (donc \uline{non vide}), et satisfait \deleted{\uline{non vide}, et satisfait la condition} (en plus de la condition \uline{PS$1'_{\mathrm g}$ de} a)) la condition \uline{PS$2_{\mathrm g}$} suivante :
% slanted
\margin{i.e. trois conditions a) $X\ne\varnothing$ b) PS$1'$ habituelle c) PS2 (double flèche)}
Pour tte double flèche
\diagram{\begin{tikzcd}[ampersand replacement=\&,column sep=small]
 a\arrow[r,shift left,"u"]\arrow[r,shift right,"\unsure{v|u}"'] \& b,
\end{tikzcd}}
existe flèche $c\xrightarrow{w}a$ telle que $\unsure{uw|wu}=vw$.
% slanted
\margin{\deleted{\unsure{$X$ 0-connexe}}}

Il revient au $=$ de dire que \unsure{pour} \deleted{$X$ soit 0-connexe et satisfait les conditions PS1 (plus faible que PS1')} \inserted{\unsure{plus succ.,|quant à ceci,}} $X$ satisfait les trois conditions a') $X$ \deleted{est} 0-connexe (plus fort que $X\ne\varnothing$), b') PS$1$ ou $X$ loc. fbt. filtr. à g \deleted{(plus faible que PS1')} et c') PS2 comme avant.

d) On dit que $X$ est \uline{loc. filtrante à g} si les composantes connexes sont \inserted{filtr. à g.}, i.e. \inserted{si $X$ satisfait les conditions PS1 et PS2 (``pseudofiltrante'' dans la terminologie SGA 4)}
% slanted
\margin{\uline{NB} Si $X$ est loc filtr. à g., il en est de $=$ de tout $X'$ étale sur $X$, car c'est évident pour PS1 qui concerne les $X_{/x}$, et pour PS2 cela résulte des \unsure{fait|faits} que $X'\to X$ est fidèle \unsure{$+$} \unsure{le fait|le fait que} \unsure{que}. \deleted{Mais} si $X$ \unsure{est} \unsure{filtr.} les $X_{/a}$ sont \unsure{filt.}, \unsure{on} \unsure{le vérifier|le vérifie}, \unsure{soit} $\unsure{f\colon}$ \illegible{1}}
On a les notions duales : \uline{faiblement filtrante à droite} (ou \unsure{simple\textsuperscript{t}|simplex} fbt. filtrante) et \uline{filtrante à droite} ou simplement filtrante. \deleted{\unsure{Les notions}} (notions duales de \textcircled{a}, \textcircled{c} ci-dessus).
Pour la notion duale de \textcircled{b}, savoir que les $X_{x\backslash}$ soient fbt. filtrantes, on ne peut plus l'appeler ``localement'' fbt filtr. \unsure{et} mais ce serait \uline{\uline{co}localement faiblement filtrante}.
\uline{Ex} $X=($
\diagram{\begin{tikzcd}[ampersand replacement=\&,column sep=small]
 0\arrow[r,"\unsure{\delta|\sigma}"] \&1\arrow[r,shift left,"\alpha'"]\arrow[r,shift right,"\beta'"'] \&2\arrow[r,shift left,"\alpha"]\arrow[r,shift right,"\beta"'] \&3
\end{tikzcd}}
$)$ avec $(\alpha\alpha'=\beta\beta')\ne(\alpha\beta'=\beta\alpha')$ \deleted{et $\alpha'\delta\ne\beta'\delta$ \unsure{comme morphismes}} \inserted{$\alpha'\delta=\beta'\delta$,}
\deleted{$X$ connexe, \unsure{soit un} $X=($}
\deleted{\diagram{\begin{tikzcd}[ampersand replacement=\&,column sep=small]
 1\arrow[r,shift left,"\alpha'"]\arrow[r,shift right,"\beta'"'] \&2\arrow[r,shift left,"\alpha"]\arrow[r,shift right,"\beta"'] \&3
\end{tikzcd}}}
\deleted{avec \unsure{$\alpha\alpha'=\beta\beta'$} \unsure{$\beta\alpha'=\beta\beta'$},} \inserted{$\alpha\ne\alpha'$, $\beta\ne\beta'$,}
